%------------------------------------------------------------------------
% Orientações Gerais
%------------------------------------------------------------------------

\chapter{ORIENTAÇÕES GERAIS}
\label{chap:orientacoes}

% ILUSTRAÇÕES
\section{Ilustrações}
\label{sec:ilustracoes}

A seguir exemplifica-se como inserir ilustrações no corpo do trabalho. As ilustrações serão indexadas automaticamente em suas respectivas listas. A numeração sequencial de figuras, tabelas e equações também ocorre de modo automático.

Referências cruzadas são obtidas por meio dos comandos \verb|\label{}| e \verb|\ref{}|. Sendo assim, não é necessário, por exemplo, saber que o número de certo capítulo é \ref{chap:orientacoes} para colocar o seu número no texto. Outra forma que pode ser utilizada é esta: \autoref{chap:orientacoes}, facilitando a inserção, remoção e manejo de elementos numerados no texto sem a necessidade de renumerar todos esses elementos.

% FIGURAS
\section{Figuras}
\label{sec:figuras}

Exemplo de como inserir uma figura. A \autoref{fig:figura-exemplo1} aparece automaticamente na lista de figuras. Os arquivos das figuras devem ser armazenados no diretório de "/Figuras". Para saber mais sobre o uso de imagens no \LaTeX{} consulte literatura especializada \cite{GOOSSENS_2007}.

\begin{figure}[!htb]
    \centering
    \caption{Exemplo de Figura}
    \includegraphics[width=0.7\textwidth]{Figuras/Figura2}
    \caption*{Fonte: Adaptado de \citeonline[p.~19]{NUNES_2014}}
    \label{fig:figura-exemplo1}
\end{figure}

% TABELAS
\section{Tabelas}
\label{sec:tabelas}

Exemplo de como inserir a \autoref{tab:tabela-exemplo1}, que aparece automaticamente na sua lista. O número da tabela e o título vem acima da tabela, e a fonte, deve vir abaixo. Para saber mais informações sobre a construção de tabelas no \LaTeX{} consulte literatura especializada \cite{MITTELBACH_2004}.

\begin{table}[!htb]
    \centering
    \caption{Tabela exemplo.}
    \label{tab:tabela-exemplo1}
    \begin{tabular}{l|rrr|c|c}
        & Valores 1 & Valores 2 & Valores 3 & Total & Total Geral \\ \hline
        Caso 1 &  86 &  77 &  81 & 244 & \multirow{3}{*}{660}\\
        Caso 2 &  19 &  74 &  25 & 118 & \\
        Caso 3 & 100 &  98 & 100 & 298 & \\
    \end{tabular}
    \vspace{0.5pt}\caption*{Fonte: Elaboração própria.}
\end{table}

% EQUAÇÕES
\section{Equações}
\label{sec:equacoes}

Exemplo de como inserir a \autoref{eq:equacao-exemplo1} e a Eq. \ref{eq:equacao-exemplo2} no corpo do texto. Observe que foram utilizadas duas formas distintas para referenciar as equações.

\begin{equation}
    X(s) = \int\limits_{t = -\infty}^{\infty} x(t) \, \text{e}^{-st} \, dt
    \label{eq:equacao-exemplo1}
\end{equation}

\begin{equation}
    F(u, v) = \sum_{m = 0}^{M - 1} \sum_{n = 0}^{N - 1} f(m, n) \exp \left[ -j 2 \pi \left( \frac{u m}{M} + \frac{v n}{N} \right) \right]
    \label{eq:equacao-exemplo2}
\end{equation}

% ALGORITMOS
\section{Algoritmos}
\label{sec:algoritmos}

Exemplo de como inserir um algoritmo. Para inserção de algoritmos utiliza-se o pacote {\ttfamily algorithm2e} que já está devidamente configurado dentro do template. 

\begin{algorithm}[ht]
    \SetKwInOut{Input}{input}
    \SetKwInOut{Output}{output}
    \SetKw{KwRetorna}{return}
    \Input{$s', l, u, step$}
    \Output{Refined solution $s''$}
    $s'' \leftarrow s'$;
    $improve \leftarrow TRUE$;
    \While{$improve$}{
        $improve \leftarrow FALSE$;
        \For{$i \in s'$}{
            Let $s_{+}'$ be the result of adding the value $step$ to position $i$ of solution $s'$\;
                \If{$l(i) \leq s_{+}'(i) \leq u(i)$}{
                    \If{$f(s_{+}') < f(s'')$}{    			   
                        $s'' \leftarrow s_{+}'$;
                        }
		      }
        }    
        \If{$f(s'') < f(s')$}{
            $s' \leftarrow s''$;
		    $improve \leftarrow TRUE$;
        }
    }
    \KwRetorna $s''$\;
    \caption{Local search procedure}
    \label{alg:busca_local}
\end{algorithm}

% LISTAS
\section{Listas}
\label{sec:lista}

Para construir listas de "\textit{bullets}"{} ou listas enumeradas, inclusive listas aninhadas, é utilizado o pacote \verb|paralist|.

\subsection{Comando \textit{itemize}}
Exemplo de duas listas não numeradas aninhadas, utilizando o comando \verb|\itemize|. Observe a indentação, bem como a mudança automática do tipo de "\textit{bullet}"{} nas listas aninhadas.

\begin{itemize}
    \item item não numerado 1
    \item item não numerado 2
    \begin{itemize}
        \item subitem não numerado 1
        \item subitem não numerado 2
        \item subitem não numerado 3
    \end{itemize}
    \item item não numerado 3
\end{itemize}

\subsection{Comando \textit{enumerate}}
Exemplo de duas listas numeradas aninhadas, utilizando o comando \verb|\enumerate|. Observe a numeração progressiva e indentação das listas aninhadas.

\begin{enumerate}
    \item item numerado 1
    \item item numerado 2
    \begin{enumerate}
        \item subitem numerado 1
        \item subitem numerado 2
        \item subitem numerado 3
    \end{enumerate}
    \item item numerado 3
\end{enumerate}

% CITAÇÕES
\section{Citações}
\label{sec:citacoes}

Citações são trechos de texto ou informações obtidas de materiais consultados quando da elaboração do trabalho. São utilizadas no texto com o propósito de esclarecer, completar e embasar as ideias do autor. Todas as publicações consultadas e utilizadas (por meio de citações) devem ser listadas, obrigatoriamente, nas referências bibliográficas, para preservar os direitos autorais. São classificadas em citações indiretas e diretas.

\subsection{Citações indiretas}
\label{subsec:citacoes_indiretas}

É a transcrição, com suas próprias palavras, das ideias de um autor, mantendo-se o sentido original. A citação indireta é a maneira que o pesquisador tem de ler, compreender e gerar conhecimento a partir do conhecimento de outros autores. Quanto à chamada da referência, ela pode ser feita de duas maneiras distintas, conforme o nome do(s) autor(es) façam parte do seu texto ou não. Exemplo de chamada fazendo parte do texto: \\
\\Enquanto \citeonline{MATURANA_2003} defendem uma epistemologia baseada na biologia. Para os autores, é necessário rever \ldots.\\

A chamada de referência foi feita com o comando \verb|\citeonline{chave}|, que produzirá a formatação correta.

A segunda forma de fazer uma chamada de referência deve ser utilizada quando se quer evitar uma interrupção na sequência do texto, o que poderia, eventualmente, prejudicar a leitura. Assim, a citação é feita e imediatamente após a obra referenciada deve ser colocada entre parênteses. Porém, neste caso específico, o nome do autor deve vir em caixa alta, seguido do ano da publicação. Exemplo de chamada não fazendo parte do texto: \\
\\Há defensores da epistemologia baseada na biologia que argumentam em favor da necessidade de \ldots \cite{MATURANA_2003}.\\

Nesse caso a chamada de referência deve ser feita com o comando \verb|\cite{chave}|, que produzirá a formatação correta.

\subsection{Citações Diretas}
\label{subsec:citacoes_diretas}

É a transcrição ou cópia de um parágrafo, de uma frase, de parte dela ou de uma expressão, usando exatamente as mesmas palavras adotadas pelo autor do trabalho consultado.

Quanto à chamada da referência, ela pode ser feita de qualquer das duas maneiras já mencionadas nas citações indiretas, conforme o nome do(s) autor(es) façam parte do texto ou não. Há duas maneiras distintas de se fazer uma citação direta, conforme o trecho citado seja longo ou curto.

Quando o trecho citado é longo (4 ou mais linhas) deve-se usar um parágrafo específico para a citação, na forma de um texto recuado (4 cm da margem esquerda), com tamanho de letra menor e espaçamento entrelinhas simples. Exemplo de citação longa:

\begin{adjustwidth}{4cm}{}
    \SingleSpacing
    \small
    Desse modo, opera-se uma ruptura decisiva entre a reflexividade filosófica, isto é a possibilidade do sujeito de pensar e de refletir, e a objetividade científica. Encontramo-nos num ponto em que o conhecimento científico está sem consciência. Sem consciência moral, sem consciência reflexiva e também subjetiva. Cada vez mais o desenvolvimento extraordinário do conhecimento científico vai tornar menos praticável a própria possibilidade de reflexão do sujeito sobre a sua pesquisa \cite[p.~28]{SILVA_2000}.
\end{adjustwidth}

Para fazer a citação longa deve-se utilizar os seguintes comandos:
\begin{verbatim}
\begin{adjustwidth}{4cm}{}
    \SingleSpacing
    \small
    <texto da citacao>
\end{adjustwidth}
\end{verbatim}

%Para fazer a citação longa deve-se utilizar os seguintes comandos:
%\begin{verbatim}
%\begin{citacao}
%<texto da citacao>
%\end{citacao}
%\end{verbatim}

No exemplo acima, para a chamada da referência o comando \verb|\cite[p.~28]{SILVA_2000}| foi utilizado, visto que os nomes dos autores não são parte do trecho citado. É necessário também indicar o número da página da obra citada que contém o trecho citado.

Quando o trecho citado é curto (3 ou menos linhas) ele deve inserido diretamente no texto entre aspas. Exemplos de citação curta: \\
\\A epistemologia baseada na biologia parte do princípio de que "assumo que não posso fazer referência a entidades independentes de mim para construir meu explicar" \cite[p.~35]{MATURANA_2003}. \\
\\A epistemologia baseada na biologia de \citeonline[p.~35]{MATURANA_2003} parte do princípio de que "assumo que não posso fazer referência a entidades independentes de mim para construir meu explicar".

\subsection{Detalhes sobre as chamadas de referências}
\label{subsec:chamadas_referencias}

Outros exemplos de comandos para as chamadas de referências e o resultado produzido por estes: \\
\\\citeonline{MATURANA_2003} \ \ \  \verb|\citeonline{MATURANA_2003}|\\
\citeonline{SILVA_2000} \ \ \   \verb|\citeonline{SILVA_2000}|\\
\cite[p.~28]{SILVA_2000} \ \ \  \verb|\cite[p.~28]{SILVA_2000}|\\
\citeonline[p.~33]{SILVA_2000} \ \ \   \verb|\citeonline[p.~33]{SILVA_2000}|\\
\cite[p.~35]{MATURANA_2003} \ \ \   \verb|\cite[p.~35]{MATURANA_2003}|\\
\citeonline[p.~35]{MATURANA_2003} \ \ \   \verb|\citeonline[p.~35]{MATURANA_2003}|\\
\cite{SILVA_2000,MATURANA_2003} \ \ \   \verb|\cite{SILVA_2000,MATURANA_2003}|\\

\subsection{Referências Bibliográficas}
\label{subsec:referencias_bibliograficas}

A bibliografia é feita no padrão \textsc{Bib}\TeX{}. As referências são colocadas em um arquivo separado. Neste template as referências são armazenadas no arquivo "referencias.bib".

Existem diversas categorias documentos e materiais componentes da bibliografia. A classe abn\TeX{} define as seguintes categorias (entradas):

\begin{verbatim}
@book
@inbook
@article
@phdthesis
@mastersthesis
@monography
@techreport
@manual
@proceedings
@inproceedings
@journalpart
@booklet
@patent
@unpublished
@misc
\end{verbatim}

Cada categoria (entrada) é formatada pelo pacote \verb|abntex2| de uma forma específica. Algumas entradas foram introduzidas especificamente para atender à norma \textbf{NBR6023:2002}, são elas: \verb|@monography|, \verb|@journalpart|,\verb|@patent|. As demais entradas são padrão \textsc{Bib}\TeX{}.

% NOTAS DE RODAPÉ
\section{Notas de Rodapé}
\label{sec:notas_rodape}

As notas de rodapé pode ser classificadas em duas categorias: notas explicativas\footnote{é o tipo mais comum de notas que destacam, explicam e/ou complementam o que foi dito no corpo do texto, como esta nota de rodapé, por exemplo.} e notas de referências. A notas de referências, como o próprio nome já indica, são utilizadas para colocar referências e/ou chamadas de referências sob certas condições.